%% title: Backing off to genus or family: thresholds per rank
%% target: Sections/Methods.tex
%% anchor: A difference between two heads is therefore claimed only when it exceeds the spread across three training seeds, not only the bootstrap interval.
%% action: after-paragraph
%% why: Defines the abstention policy the results use (sec:results_backoff): one threshold per rank, each fitted on the images that reach that rank, and the expected-utility rule it is compared with.
%% source: journal/2026-10-23-what-is-the-best-back-off-policy.md (Q: policies P0-P3, eleven splits); paper/DRAFT.md §4.6 (C1: conditional fitting, genus precision 0.487 vs 0.970); dev/083_hf_release.py (cascade); dev/097_backoff_policies.py
%% check: Paste below step 11. | The Deng et al. (2012) reference is from memory; verify. | Temperatures (1.91, 1.70) belong to the recipe-L release models; the rerun's models get their own.
%% status: ready
\subsection{Backing off to a coarser rank}\label{sec:backoff}
Summing the species posterior over each genus and family gives a posterior at every rank from one forward pass, so a model can answer at a coarser rank without a second model.
The policy is greedy: answer at species if the top species confidence reaches a threshold $\tau_s$; otherwise at genus if the top genus confidence reaches $\tau_g$; otherwise at family if it reaches $\tau_f$; otherwise abstain.
Each threshold is the lowest that gives 95\,\% precision on the calibration images that reach its rank, that is, the images that failed every finer rank.
Fitting the coarse thresholds on all images instead lets them borrow reliability from easy images: on the 15,453 test images whose species confidence is below 0.40, genus precision is 0.487, against 0.970 over all images.

On the trap images the thresholds are fitted on half of the (trap, night) groups and evaluated on the other half, over eleven random halvings.
Heads trained with an angular margin saturate their softmax (B8 returns a confidence of exactly 1.0 for 70\,\% of trap images), so they are first rescaled by a temperature fitted on in-distribution validation data.
The greedy policy is compared with an expected-utility rule, which answers at the rank $r$ that maximises $u_r p_r - c\,(1-p_r)$, where $p_r$ is the top posterior at rank $r$, and abstains when every rank gives a negative value \textbf{(REF: Deng et al. 2012, hedging your bets)}.
The utilities are 1, 0.8 and 0.5 for species, genus and family, and the error cost $c$ is chosen to reach the same 95\,\% precision overall.
Each policy is reported as the share of images answered at each rank, the precision of each returned rank, and the share answered correctly at some rank.
